\subsection{The Big Data Deluge}\label{subsec:big-data--the-big-data-deluge}

The phrase \definition{Big Data deluge} describes the idea that \hl{modern systems are producing data at a scale and rate that traditional data-management approaches were not designed for.} The key idea is that \hl{data generation is growing faster than our ability to manage it traditionally.}

\highspace
The word \emph{deluge} is deliberate: it suggests a flood. We are not just dealing with ``\emph{more rows in a database}''. We are dealing with a \hl{continuous accumulation of data from a huge number of sources.}

\highspace
\begin{flushleft}
    \textcolor{Green3}{\faIcon{project-diagram} \textbf{Where does all this data come from?}}
\end{flushleft}
Historically, most digital data came from relatively controlled business systems such as: transactions, customer records, inventory, banking operations, and enterprise applications. These sources were usually \hl{structured and generated at a relatively predictable rate.}

\highspace
Today, data comes from many more places:
\begin{equation*}
    \text{Data Sources} = \text{Humans} + \text{Software} + \text{Machines} + \text{Sensors}
\end{equation*}
For example smartphone interactions, web clicks, GPS positions, social-media activity, server logs, IoT sensors, cameras, industrial machinery, payment systems, connected vehicles, and wearable devices. So the number of \textbf{data-\break producing entities} has increased dramatically.

\highspace
\begin{flushleft}
    \textcolor{Green3}{\faIcon{cogs} \textbf{Machines now generate data automatically}}
\end{flushleft}
Traditional data was often created because a human explicitly performed a business action. For example:
\begin{equation*}
    \text{customer buys product}
    \rightarrow
    \text{transaction record}
\end{equation*}
Modern systems now \hl{generate much more data around the same action}:  user opened an app, searched for a product, scrolled the page, clicked an image, added the product to the cart, removed it, added it again, paid, recorded payment latency, updated delivery location, and opened a notification.

\highspace
A \hl{single business event can now produce tens or hundreds of technical and behavioural events.} Therefore:
\begin{equation*}
    \text{one real-world action}
    \rightarrow
    \text{many digital observations}
\end{equation*}
This is \hl{one reason why data volume has exploded.}

\highspace
\begin{flushleft}
    \textcolor{Green3}{\faIcon{database} \textbf{Storage is no longer only about final business records}}
\end{flushleft}
In a traditional information system, we might mainly keep the final result:
\begin{lstlisting}[mathescape=true]
Order #1234
Customer = Alice
Total = $\texttt{\euro}$ 84
\end{lstlisting}
But \hl{modern analytics often values the entire process that led to that result.} For example search, click, view, pause, scroll, add-to-cart, remove-from-cart, purchase.

\highspace
\textcolor{Green3}{\faIcon{question-circle} \textbf{Why?}} Because those intermediate events reveal \textbf{behaviour}. The final transaction tells us: ``\emph{Alice bought product X}''. But the complete event stream tells us much more: ``\emph{Alice searched three alternatives, compared prices, hesitated for five minutes, abandoned the cart, returned later, and then completed the purchase.}'' That additional information can be far more useful for prediction and personalization.

\highspace
\begin{flushleft}
    \textcolor{Green3}{\faIcon{chart-area} \textbf{From scarce data to abundant data}}
\end{flushleft}
Historically, one difficulty was: ``\emph{\textbf{how can we obtain enough data?}}''. Now, in many contexts, the problem has become: ``\emph{\textbf{how can we manage all the data we already produce?}}''. This is a fundamental shift.

\highspace
Previously, collecting data was expensive. Therefore organizations tended to decide carefully: ``\emph{which data should we store?}''. \hl{Today, storage is relatively cheaper and automatic instrumentation is common, so the instinct often becomes: ``\emph{store as much as possible because it may become useful later}''.} This produces a huge expansion in stored information.

\highspace
\begin{flushleft}
    \textcolor{Green3}{\faIcon{expand-arrows-alt} \textbf{The deluge is not only about Volume}}
\end{flushleft}
Although \emph{deluge} immediately makes us think of \textbf{quantity}, the problem is broader. Imagine a smart city. It may collect traffic sensor readings, bus positions, weather measurements, CCTV video, social media, mobile phone signals, public-event data, parking information, and pollution sensors. The \textbf{challenge} is not merely that there are many bytes. The \textbf{data}:
\begin{itemize}
    \item \textbf{Differs} in \textbf{type} (structured, semi-structured, unstructured);
    \item \textbf{Arrives} at different \textbf{rates} (from milliseconds to hours);
    \item \textbf{Has} different levels of \textbf{quality} (from highly accurate to noisy or missing).
\end{itemize}
This anticipates the Big Data dimensions we will later formalize as: Volume, Variety, Velocity, and Veracity. So the ``\emph{deluge}'' is the phenomenon; the \textbf{Vs} (i.e., the four Vs) later give us a \textbf{framework for describing its properties}.

\highspace
\begin{flushleft}
    \textcolor{Red2}{\faIcon{exclamation-triangle} \textbf{Why traditional systems start struggling}}
\end{flushleft}
Suppose a company has a traditional architecture: one Relational DataBase Management System (DBMS). This may work perfectly well for a certain workload. But suppose the company begins ingesting $10^3$ events per second. Then $10^5$ and eventually $10^7$ events per second. At the same time, it stores text, JSON, images, logs, user interactions, and sensor streams.

\highspace
At some point, simply \hl{making one server larger becomes increasingly difficult or expensive.} This creates pressure toward distributed storage and computation. Instead of one very large computer (\textbf{scale up}), we can distribute the work across many machines (\textbf{scale out}). This will later connect directly to horizontal scalability, Hadoop, distributed databases, sharding, and replication.

\highspace
\begin{flushleft}
    \textcolor{Green3}{\faIcon{filter} \textbf{Data generation vs. useful information}}
\end{flushleft}
One subtle point is that:
\begin{equation*}
    \text{more data} \neq \text{more knowledge}
\end{equation*}
A company might collect 1 PB of data and still learn almost nothing from it. The \hl{value appears only if the organization can transform that data.} This gives us the pipeline:
\begin{equation*}
    \text{Data}
    \rightarrow
    \text{Processing}
    \rightarrow
    \text{Information}
    \rightarrow
    \text{Insight}
    \rightarrow
    \text{Action}
\end{equation*}
The deluge creates both: \textbf{a problem} because we need to manage enormous quantities of data, and \textbf{an opportunity} because that data potentially contains useful information.

\highspace
\begin{examplebox}[: The Big Data deluge in practice]
    Consider a streaming platform. A very simple historical dataset might have been:
    \begin{center}
        \texttt{user\_id}, \texttt{movie\_id}, \texttt{rating}, \texttt{date}
    \end{center}
    Now suppose the platform records:
    \begin{itemize}
        \item When playback starts
        \item When playback stops
        \item Pause events
        \item Rewind events
        \item Device type
        \item Network quality
        \item Search queries
        \item Time spent browsing
        \item Which thumbnail was shown
        \item Whether the thumbnail was clicked
        \item Percentage of episode watched
    \end{itemize}
    Instead of one record per movie interaction, we might now have hundreds of events. For one user, that might be 100 events per day. For $10^8$ users, that is $100 \times 10^8 = 10^{10}$ events per day. \hl{That simple multiplication explains why distributed data architectures become necessary very quickly.}
\end{examplebox}

\highspace
\begin{flushleft}
    \textcolor{Green3}{\faIcon{sync-alt} \textbf{The feedback effect}}
\end{flushleft}
Another important consequence is that successful digital applications often generate \textbf{even more data as they grow}. Suppose a service becomes popular, then:
\begin{enumerate}
    \item More users use the service
    \item Which produces more interactions
    \item Which produces more data
    \item Which may enable better analytics (assuming a good data pipeline)
    \item Which improves the service and product
    \item Which attracts more users
\end{enumerate}
This already starts to look like a \textbf{virtuous cycle}, which is exactly what the next section will formalize.

\highspace
\begin{takeawaysbox}
    The \textbf{Big Data deluge} means that digital systems generate and retain unprecedented amounts of data from a rapidly increasing number of sources.

    \highspace
    It produces three important consequences.
    \begin{enumerate}
        \item Traditional centralized architectures may stop scaling.
        \item We need new ways to store and process heterogeneous, fast-growing data.
        \item The large amount of data creates opportunities for better analysis and decisions.
    \end{enumerate}
    So the deluge is both a \textbf{problem} and an \textbf{opportunity}:
    \begin{equation*}
        \text{Data Deluge}
        \rightarrow
        \text{Need process data}
        \rightarrow
        \text{Extract insights}
        \rightarrow
        \text{Create value}
    \end{equation*}
\end{takeawaysbox}
